\begin{abstract}
We investigate the security implications of Apple's data memory-dependent prefetcher (DMP), beginning with the M1. Our scope has broadened from reproducing GoFetch to studying DMP exploitation through both Augury and GoFetch. At the midpoint, our experimental work has focused on Augury activation and reverse engineering. We report an Augury existence sweep across 14 training sizes and 81 offsets per size on an Apple M1. The supplied sweep summary reports activation at 12 training sizes, with a peak median speedup of 16.3$\times$ among activating offsets at 256 training pointers. We describe the measurement limits, an unresolved non-activating case at eight training pointers, and the remaining exploitation demonstrations. Our planned novel investigation compares selected activation conditions under Asahi Linux on M1 and on an M4 chip, with particular attention to distinguishing changed behavior from evidence of mitigation.
\end{abstract}
